\documentclass[10pt,letterpaper]{article}
\usepackage{cornell-notes}

\title{Chapter 26: The M\"obius Function}
\author{}
\date{}

\setDocTitle{Chapter 26: The M\"obius Function}
\setDocSubtitle{Combinatorial Algorithms}
\setDocVersion{v1.0}
\setDocStatus{Study Companion}
\setDocOwner{Combinatorial Algorithms Cornell Notes}
\setDocAuthor{}
\setDocDate{\today}
\setDocSummary{Cornell-format study and review notes for Chapter 26: The M\"obius Function.}

\setCornellCollection{Combinatorial Algorithms Cornell Notes}
\setCornellUnitType{Chapter}
\setCornellUnitNumber{26}
\setCornellUnitTitle{The M\"obius Function}
\setCornellCompanionLabel{Study and review companion}

\hypersetup{pdftitle={Chapter 26: The M\"obius Function},pdfsubject={Cornell notes},pdfauthor={}}

\begin{document}
\maketitle
\begin{tcolorbox}[colback=Paper,colframe=Ink]{\small\bfseries\color{Teal}CORNELL NOTES}\par{\LARGE\bfseries Chapter 26: The M\"obius Function (MOBIUS)}\par\medskip\textbf{Topic:} Poset incidence algebra by triangular matrix operations\hfill\textbf{Date:} \today\par\textbf{Study purpose:} Compute the M\"obius matrix from a finite order relation.
\end{tcolorbox}
\begin{tcolorbox}[colback=Teal!5,colframe=Teal,title=Essential question]How can covering relations, triangular numbering, and two inversions recover the M\"obius function of a finite poset?\end{tcolorbox}
\begin{tcolorbox}[colback=white,colframe=Mist,title=Learning objectives]\begin{itemize}\item Relate the zeta and M\"obius matrices.\item Recover reachability from a cover matrix.\item Trace the MOBIUS pipeline and relabeling steps.\item Justify the use of triangular inversion.\item State the principal time and storage costs.\end{itemize}\textbf{Prerequisites:} finite posets, cover relations, matrix inverses, linear extensions.\end{tcolorbox}
\section*{Cornell notes}
\begin{longtable}{@{}Q!{\color{Mist}\vrule width 1pt}A@{}}\cornellhead
\cue{What are $Z$ and $M$?}{For a finite poset, the zeta matrix satisfies $Z_{ij}=1$ when $i\leq j$ and $0$ otherwise. Its inverse $M=Z^{-1}$ contains the M\"obius values. Thus $\sum_{i\leq k\leq j}\mu(i,k)=\delta_{ij}$.}
\cue{Why start from covers?}{Let $H$ be the covering-relation matrix in a consistent order. The entry $(H^r)_{ij}$ counts saturated chains of length $r$ from $i$ to $j$. Because $H$ is strictly upper triangular, it is nilpotent and $(I-H)^{-1}=I+H+H^2+\cdots$ is finite.}
\cue{How is zeta recovered?}{An entry of $(I-H)^{-1}$ is nonzero exactly when a chain connects the two elements, including the length-zero chain on the diagonal. Replacing every nonzero entry by $1$ therefore yields $Z$. Chain multiplicities are deliberately discarded.}
\cue{MOBIUS pipeline}{\begin{enumerate}\item Use TRIANG to find a consistent numbering.\item Apply RENUMB to put $H$ in triangular form.\item Invert $I-H$.\item Threshold nonzero entries to obtain $Z$.\item Invert $Z$ to obtain $M$.\item Invert the label permutation and renumber rows and columns back.\end{enumerate}}
\cue{Key invariant}{Every matrix inverted is unit upper triangular. Consequently its inverse exists over the integers, remains upper triangular, and can be computed by the specialized recurrence from the preceding matrix chapter.}
\cue{Why relabel twice?}{Triangular inversion requires an order compatible with the poset, but the caller expects answers in the original labels. The forward permutation enables efficient arithmetic; the inverse permutation restores the semantic identity of every element.}
\cue{Complexity and storage}{The dense triangular products and inversions require $O(n^3)$ arithmetic in the stated representation; renumbering is $O(n^2)$. The operations can reuse a single integer matrix plus permutation work arrays, although intermediate entries can exceed $0$ and $1$.}
\cue{Failure checks}{The input must represent an acyclic partial order. A diagonal cover, a directed cycle, or inconsistent relabeling invalidates nilpotence and the triangular assumptions. Thresholding must occur only after the first inverse, not after the final M\"obius inverse.}
\cue{Retrieval practice}{\begin{enumerate}\item Why is $H$ nilpotent after triangular numbering?\item What information is lost by thresholding?\item Why are two inversions needed?\item Which permutation restores the original labels?\item What identity verifies the output?\end{enumerate}}
\end{longtable}
\begin{tcolorbox}[colback=Ink!4,colframe=Ink,title=Cornell summary]
The M\"obius matrix is the inverse of the zeta matrix. MOBIUS first converts a cover relation to a triangular representation. The finite series $(I-H)^{-1}$ counts chains, whose nonzero pattern is precisely the zeta relation. Thresholding that matrix and inverting once more produces the M\"obius values. A final inverse renumbering returns those values to the original element labels. The method is algebraically clean and storage-efficient for dense matrix input, with cubic arithmetic cost dominated by triangular matrix work.
\end{tcolorbox}
\end{document}
